% GENERATED FILE. Edit canonical lessons, then run npm run generate:books.
% canonical-source-hash: fnv1a64:a5308c1ff87da8be
% canonical-lessons: GU-C180-mulakat, GU-C180-amantran, GU-C180-rah-jovi, GU-C180-rokavu, GU-C180-vitavvu

\chapter{Visits and invitations}
\label{ch:gu-visits-and-invitations}
\hlchaptermodality{180}

\begin{chapteropening}
\textbf{By the end of this chapter:} \emph{I can talk about invitations, visits and guests who will stay.}

Complete the last lesson of chapter 180: I can talk about invitations, visits and guests who will stay.
\end{chapteropening}

\section[\texorpdfstring{\gu{મુલાકાત} (mulākāt)}{મુલાકાત (mulākāt)}]{\gu{મુલાકાત} (mulākāt) — visit, meeting}

\label{lesson:GU-C180-mulakat}



\begin{quote}
Before the new one: say the Gujarati for beforehand, then the Gujarati for certain.
\end{quote}

\subsection*{You'll want to know: \gu{મુલાકાત}}
\textbf{\gu{મુલાકાત}} — \emph{mulākāt} — \textquotedblleft{}visit, meeting\textquotedblright{}.

\begin{quote}
\textbf{\gu{આવતીકાલે} \gu{અમારી} \gu{મુલાકાત} \gu{થશે}.} — \emph{āvtīkāle amārī mulākāt thaśe} — \textquotedblleft{}We will meet tomorrow.\textquotedblright{}
\end{quote}

\begin{cousinweb}
\textbf{\gu{મુલાકાત}} (\emph{mulākāt}) is a loan from Arabic \emph{mulāqāt}, \textquotedblleft{}meeting\textquotedblright{}, which reached Gujarati through Persian. Its Arabic root is \emph{l-q-y}, \textquotedblleft{}to meet\textquotedblright{} (Lane).
\end{cousinweb}

\begin{grammarlens}[title={What you've built}]
One more word for everyday life at A2.
\end{grammarlens}

\subsection*{Your turn}
\begin{itemize}
\raggedright
  \item \emph{Say it:} \emph{mulākāt}
  \item \emph{Say it:} \emph{mulākāt}, once more
  \item \emph{Say it:} \emph{āvtīkāle amārī mulākāt thaśe} — We will meet tomorrow.
  \item \emph{Recall:} say the Gujarati for beforehand, then the Gujarati for certain, then say \emph{mulākāt} again
\end{itemize}

\begin{tcolorbox}[breakable,colback=teal!4,colframe=teal!35!black,title={Before you move on}]
What does \gu{મુલાકાત} mean? (\textquotedblleft{}Visit, meeting\textquotedblright{}.) Say it once more.
\end{tcolorbox}
\section[\texorpdfstring{\gu{આમંત્રણ} (āmantraṇ)}{આમંત્રણ (āmantraṇ)}]{\gu{આમંત્રણ} (āmantraṇ) — invitation}

\label{lesson:GU-C180-amantran}



\begin{quote}
Before the new one: say the Gujarati for certain, then the Gujarati for visit.
\end{quote}

\subsection*{You'll want to know: \gu{આમંત્રણ}}
\textbf{\gu{આમંત્રણ}} — \emph{āmantraṇ} — \textquotedblleft{}invitation\textquotedblright{}.

\begin{quote}
\textbf{\gu{હું} \gu{તેમને} \gu{આમંત્રણ} \gu{આપીશ}.} — \emph{hũ temne āmantraṇ āpīś} — \textquotedblleft{}I will invite them.\textquotedblright{}
\end{quote}

\begin{cousinweb}
\textbf{\gu{આમંત્રણ}} (\emph{āmantraṇ}) is Sanskrit \emph{āmantraṇa}, \textquotedblleft{}addressing, calling; invitation\textquotedblright{}, taken over whole as a learned loan, a tatsama.
\end{cousinweb}

\begin{grammarlens}[title={What you've built}]
One more word for everyday life at A2.
\end{grammarlens}

\subsection*{Your turn}
\begin{itemize}
\raggedright
  \item \emph{Say it:} \emph{āmantraṇ}
  \item \emph{Say it:} \emph{āmantraṇ}, once more
  \item \emph{Say it:} \emph{hũ temne āmantraṇ āpīś} — I will invite them.
  \item \emph{Recall:} say the Gujarati for certain, then the Gujarati for visit, then say \emph{āmantraṇ} again
\end{itemize}

\begin{tcolorbox}[breakable,colback=teal!4,colframe=teal!35!black,title={Before you move on}]
What does \gu{આમંત્રણ} mean? (\textquotedblleft{}Invitation\textquotedblright{}.) Say it once more.
\end{tcolorbox}
\section[\texorpdfstring{\gu{રાહ} \gu{જોવી} (rāh jovī)}{રાહ જોવી (rāh jovī)}]{\gu{રાહ} \gu{જોવી} (rāh jovī) — to wait}

\label{lesson:GU-C180-rah-jovi}



\begin{quote}
Before the new one: say the Gujarati for visit, then the Gujarati for invitation.
\end{quote}

\subsection*{You'll want to know: \gu{રાહ} \gu{જોવી}}
\textbf{\gu{રાહ} \gu{જોવી}} — \emph{rāh jovī} — \textquotedblleft{}to wait\textquotedblright{}.

\begin{quote}
\textbf{\gu{હું} \gu{સ્ટેશને} \gu{તમારી} \gu{રાહ} \gu{જોઈશ}.} — \emph{hũ sṭeśane tamārī rāh joīś} — \textquotedblleft{}I will wait for you at the station.\textquotedblright{}
\end{quote}

Word for word, to watch the road: \textbf{\gu{રાહ}}, \emph{rāh}, and \textbf{\gu{જોવું}}, \emph{jovũ}.

\begin{cousinweb}
\textbf{\gu{રાહ} \gu{જોવી}} (\emph{rāh jovī}) is built from \emph{rāh}, \textquotedblleft{}road, way\textquotedblright{}, usually traced to Persian \emph{rāh}, \textquotedblleft{}road\textquotedblright{}, and the inherited verb \emph{jovũ}, \textquotedblleft{}to see\textquotedblright{}.
\end{cousinweb}

\begin{grammarlens}[title={What you've built}]
One more word for everyday life at A2.
\end{grammarlens}

\subsection*{Your turn}
\begin{itemize}
\raggedright
  \item \emph{Say it:} \emph{rāh jovī}
  \item \emph{Say it:} \emph{rāh jovī}, once more
  \item \emph{Say it:} \emph{hũ sṭeśane tamārī rāh joīś} — I will wait for you at the station.
  \item \emph{Recall:} say the Gujarati for visit, then the Gujarati for invitation, then say \emph{rāh jovī} again
\end{itemize}

\begin{tcolorbox}[breakable,colback=teal!4,colframe=teal!35!black,title={Before you move on}]
What does \gu{રાહ} \gu{જોવી} mean? (\textquotedblleft{}To wait\textquotedblright{}.) Say it once more.
\end{tcolorbox}
\section[\texorpdfstring{\gu{રોકાવું} (rokāvũ)}{રોકાવું (rokāvũ)}]{\gu{રોકાવું} (rokāvũ) — to stay on, to stop over}

\label{lesson:GU-C180-rokavu}



\begin{quote}
Before the new one: say the Gujarati for invitation, then the Gujarati for to wait.
\end{quote}

\subsection*{You'll want to know: \gu{રોકાવું}}
\textbf{\gu{રોકાવું}} — \emph{rokāvũ} — \textquotedblleft{}to stay on, to stop over\textquotedblright{}.

\begin{quote}
\textbf{\gu{મહેમાન} \gu{બે} \gu{દિવસ} \gu{રોકાશે}.} — \emph{mahemān be divas rokāśe} — \textquotedblleft{}The guests will stay two days.\textquotedblright{}
\end{quote}

\begin{cousinweb}
\textbf{\gu{રોકાવું}} (\emph{rokāvũ}) is built on the inherited \emph{rokvũ}, \textquotedblleft{}to stop, to hold back\textquotedblright{}, with the intransitive ending \emph{-āvũ}. CDIAL traces \emph{rokvũ} to \emph{*rōkk}, \textquotedblleft{}stop\textquotedblright{}, a reconstructed form (CDIAL 10827).
\end{cousinweb}

\begin{grammarlens}[title={What you've built}]
One more word for everyday life at A2.
\end{grammarlens}

\subsection*{Your turn}
\begin{itemize}
\raggedright
  \item \emph{Say it:} \emph{rokāvũ}
  \item \emph{Say it:} \emph{rokāvũ}, once more
  \item \emph{Say it:} \emph{mahemān be divas rokāśe} — The guests will stay two days.
  \item \emph{Recall:} say the Gujarati for invitation, then the Gujarati for to wait, then say \emph{rokāvũ} again
\end{itemize}

\begin{tcolorbox}[breakable,colback=teal!4,colframe=teal!35!black,title={Before you move on}]
What does \gu{રોકાવું} mean? (\textquotedblleft{}To stay on, to stop over\textquotedblright{}.) Say it once more.
\end{tcolorbox}
\section[\texorpdfstring{\gu{વિતાવવું} (vitāvvũ)}{વિતાવવું (vitāvvũ)}]{\gu{વિતાવવું} (vitāvvũ) — to pass (time), to spend}

\label{lesson:GU-C180-vitavvu}



\begin{quote}
Before the new one: say the Gujarati for to wait, then the Gujarati for to stay on.
\end{quote}

\subsection*{You'll want to know: \gu{વિતાવવું}}
\textbf{\gu{વિતાવવું}} — \emph{vitāvvũ} — \textquotedblleft{}to pass (time), to spend\textquotedblright{}.

\begin{quote}
\textbf{\gu{અમે} \gu{રજા} \gu{ગામમાં} \gu{વિતાવીશું}.} — \emph{ame rajā gāmmā̃ vitāvīśũ} — \textquotedblleft{}We will spend the holiday in the village.\textquotedblright{}
\end{quote}

The causative of \textbf{\gu{વીતવું}}, to pass.

\begin{cousinweb}
\textbf{\gu{વિતાવવું}} (\emph{vitāvvũ}) is the causative of the inherited \emph{vītvũ}, \textquotedblleft{}to pass, of time\textquotedblright{}, which CDIAL traces to Sanskrit \emph{vṛtta}, \textquotedblleft{}passed, elapsed (of time)\textquotedblright{} (CDIAL 12069).
\end{cousinweb}

\begin{grammarlens}[title={What you've built}]
One more word for everyday life at A2.
\end{grammarlens}

\subsection*{Your turn}
\begin{itemize}
\raggedright
  \item \emph{Say it:} \emph{vitāvvũ}
  \item \emph{Say it:} \emph{vitāvvũ}, once more
  \item \emph{Say it:} \emph{ame rajā gāmmā̃ vitāvīśũ} — We will spend the holiday in the village.
  \item \emph{Recall:} say the Gujarati for to wait, then the Gujarati for to stay on, then say \emph{vitāvvũ} again
\end{itemize}

\begin{tcolorbox}[breakable,colback=teal!4,colframe=teal!35!black,title={Before you move on}]
What does \gu{વિતાવવું} mean? (\textquotedblleft{}To pass (time), to spend\textquotedblright{}.) Say it once more.
\end{tcolorbox}
